\documentclass[11pt]{article}
\usepackage[margin=1in]{geometry}
\usepackage{amsmath, amssymb}
\usepackage{graphicx}
\usepackage{booktabs}
\usepackage{hyperref}
\usepackage{bm}
\usepackage{enumitem}
\title{Taxonomy-Aware and Hyperbolic Variational Autoencoders for Microbiome Tables}
\author{biomevae Project}
\date{\today}
\begin{document}
\maketitle
\begin{abstract}
We document the methodology implemented in the \texttt{biomevae} repository for learning
low-dimensional representations of sparse microbiome abundance tables.
The system extends classical variational autoencoders (VAEs) in two complementary
ways: (i) a taxonomy-aware decoder regularized by hierarchical aggregations and graph
Laplacian smoothing, and (ii) a hyperbolic latent space realized on the Poincar\'e ball.
We derive the optimization objectives, detail the data processing pipeline, and explain
the training procedure shared across Euclidean and hyperbolic variants. This report is
intended as a rigorous companion to the command-line interface provided by the codebase.
\end{abstract}
\section{Notation and Data Model}
Let $\mathcal{F}$ denote the set of microbial features (species-level genome bins) and
$\mathcal{S}$ the set of samples. The raw input is a tab-separated matrix where rows
correspond to features and columns to samples, with metadata columns for the feature
identifier and NCBI taxonomy ID. Parsing is implemented by \texttt{load\_matrix}\footnote{\texttt{src/biomevae/data.py}}:
for an input matrix $X^{\mathrm{feat}\times \mathrm{sample}}$ the tool transposes to
\begin{equation}
    \bm{X} \in \mathbb{R}^{n \times d}, \qquad n = |\mathcal{S}|,~ d = |\mathcal{F}|,
\end{equation}
optionally applying the element-wise transform $x \mapsto \log(1+x)$.
When standardization is requested, feature-wise statistics are computed on the
training subset only, yielding mean vector $\bm{\mu}$ and standard deviation
$\bm{\sigma}$; the standardized observations are $(\bm{x}-\bm{\mu})/\bm{\sigma}$ and
$(\bm{\mu},\bm{\sigma})$ are persisted for downstream reuse.
\section{Base Variational Autoencoder}
\subsection{Encoder--Decoder Architecture}
For an input dimension $d$ and latent dimension $\ell$, the Euclidean VAE uses a
stacked multi-layer perceptron encoder $f_{\phi} : \mathbb{R}^d \to \mathbb{R}^{2\ell}$
with hidden widths $h_1,\dots,h_L$ and activation $\sigma$. Implementation details are
captured by \texttt{VAE}\footnote{\texttt{src/biomevae/models/vae.py}}:
\begin{align}
    \bm{h}^{(0)} &= \bm{x},\\
    \bm{h}^{(\ell)} &= \sigma\bigl(\mathbf{W}^{(\ell)} \bm{h}^{(\ell-1)} + \bm{b}^{(\ell)}\bigr), \qquad \ell=1,\dots,L,
\end{align}
optionally followed by layer normalization and dropout. The final hidden state produces
mean and log-variance parameters $\bm{\mu} = \mathbf{W}_{\mu} \bm{h}^{(L)} + \bm{b}_{\mu}$,
$\bm{\log\sigma^2} = \mathbf{W}_{\log} \bm{h}^{(L)} + \bm{b}_{\log}$. The decoder
mirrors the encoder with weights $\{\tilde{\mathbf{W}}^{(\ell)}\}$, terminating in a
linear reconstruction layer $g_{\theta} : \mathbb{R}^{\ell} \to \mathbb{R}^{d}$.
\subsection{Stochastic Latent Sampling}
Given $(\bm{\mu}, \bm{\log\sigma^2})$, the reparameterization trick draws
\begin{equation}
    \bm{z} = \bm{\mu} + \exp\bigl(\tfrac{1}{2}\bm{\log\sigma^2}\bigr) \odot \bm{\epsilon},
    \qquad \bm{\epsilon} \sim \mathcal{N}(\bm{0}, \mathbf{I}).
\end{equation}
A forward pass returns the reconstruction $\hat{\bm{x}} = g_{\theta}(\bm{z})$
accompanied by $(\bm{\mu},\bm{\log\sigma^2})$ for the loss computation.
\subsection{Objective Functions}
The reconstruction loss supports mean-squared error (MSE), mean absolute error (MAE),
and Huber penalties via \texttt{reconstruction\_loss}\footnote{\texttt{src/biomevae/losses.py}}. For
$N$ observations the default MAE objective reads
\begin{equation}
    \mathcal{L}_{\mathrm{rec}} = \frac{1}{N} \sum_{i=1}^N \frac{1}{d} \lVert \hat{\bm{x}}_i - \bm{x}_i \rVert_1.
\end{equation}
The KL divergence between the approximate posterior and the isotropic Gaussian prior
is evaluated per sample as
\begin{equation}
    \mathcal{D}_{\mathrm{KL}}(q_{\phi}(\bm{z}|\bm{x}) \Vert p(\bm{z}))
    = -\tfrac{1}{2} \sum_{j=1}^{\ell} \bigl(1 + \log\sigma^2_{ij} - \mu_{ij}^2 - \exp(\log\sigma^2_{ij})\bigr).
\end{equation}
The trainer combines the two terms using either the $\beta$-VAE or the capacity
objective (\texttt{compute\_losses}); setting $\beta=1$ recovers the
traditional VAE objective without annealing. A dedicated CLI,
\texttt{biomevae-train-vanilla}, exposes this configuration alongside the
existing \texttt{biomevae-train} entry point:
\begin{align}
    \mathcal{L}_{\beta\text{-VAE}} &= \mathcal{L}_{\mathrm{rec}} + \beta \cdot \mathcal{D}_{\mathrm{KL}},\\
    \mathcal{L}_{\mathrm{capacity}} &= \mathcal{L}_{\mathrm{rec}} + \gamma \bigl\lvert \mathcal{D}_{\mathrm{KL}} - C\bigr\rvert,
\end{align}
where $\beta$ is annealed according to \texttt{beta\_schedule} and $C$ follows
\texttt{capacity\_schedule}. Optional ``free bits'' enforce per-dimension lower bounds
on the KL contribution to mitigate posterior collapse.
\section{Taxonomy-Aware Regularization}
Let $\mathcal{T}$ denote the taxonomy graph containing edges between features and their
coarser ancestors (species, genus, family, etc.). The CLI \texttt{biomevae-train-tax}
parses a taxonomy table and constructs two structures:
\begin{enumerate}[label=\roman*), leftmargin=*]
    \item Aggregation matrices $\mathbf{A}^{(\ell)} \in \{0,1\}^{G_{\ell} \times d}$ mapping features to
    higher-level groups at taxonomic level $\ell$, implemented in
    \texttt{build\_taxonomy\_structures}.
    \item A feature--feature affinity matrix $\mathbf{W}$ whose entries accumulate
    positive weights when two features share the same species, genus, or family label.
    The Laplacian is $\mathbf{L} = \mathbf{D} - \mathbf{W}$ with $\mathbf{D}$ the degree
    matrix.
\end{enumerate}
\subsection{Hierarchical Reconstruction Loss}
For each selected taxonomic level $\ell$, the trainer forms aggregated observations
\begin{equation}
    \tilde{\bm{x}}^{(\ell)} = \bm{x} \mathbf{A}^{(\ell)\top}, \qquad
    \hat{\tilde{\bm{x}}}^{(\ell)} = \hat{\bm{x}} \mathbf{A}^{(\ell)\top}.
\end{equation}
The same reconstruction loss is applied to the aggregated counts, yielding
$\mathcal{L}^{(\ell)}_{\mathrm{rec}}$. The total objective becomes
\begin{equation}
    \mathcal{L} = \mathcal{L}_{\beta\text{-VAE}} + \lambda_{\mathrm{tax}} \sum_{\ell \in \mathcal{L}} \mathcal{L}^{(\ell)}_{\mathrm{rec}},
\end{equation}
where $\lambda_{\mathrm{tax}}$ is user-controlled via \texttt{--tax-loss-weight}.
This encourages reconstructions that preserve coarse-scale abundances consistent with
the provided phylogenetic hierarchy.
\subsection{Graph Laplacian Smoothing}
When enabled, the decoder output is further regularized using the quadratic form of
$\mathbf{L}$. For a minibatch matrix $\hat{\mathbf{X}} \in \mathbb{R}^{B \times d}$,
\begin{equation}
    \mathcal{S} = \frac{1}{B} \sum_{i=1}^B \hat{\bm{x}}_i^{\top} \mathbf{L} \hat{\bm{x}}_i
    = \frac{1}{B} \mathrm{tr}\bigl(\hat{\mathbf{X}} \mathbf{L} \hat{\mathbf{X}}^{\top}\bigr).
\end{equation}
The loss $\mathcal{S}$ penalizes large discrepancies between features with strong
phylogenetic affinity, acting as a smoothness prior. The weight \texttt{--tax-laplacian}
scales this term.
\section{Hyperbolic Latent Geometry}
To capture hierarchical relationships in latent space directly, the hyperbolic variant
\texttt{HyperbolicVAE}\footnote{\texttt{src/biomevae/models/hyperbolic.py}} models latent
codes on the Poincar\'e ball of curvature $-c$ with $c>0$. The encoder produces tangent
vectors $\bm{v} \in T_{\bm{0}}\mathcal{B}^d$, parameterized by $(\bm{\mu},\bm{\log\sigma^2})$ in the
Euclidean tangent space. Sampling proceeds by
\begin{equation}
    \bm{v} = \bm{\mu} + \exp\bigl(\tfrac{1}{2}\bm{\log\sigma^2}\bigr) \odot \bm{\epsilon},
    \quad \bm{\epsilon} \sim \mathcal{N}(\bm{0}, \mathbf{I}),
    \qquad \bm{z} = \mathrm{exp}_{\bm{0}}^{\mathcal{B}^d}(\bm{v}),
\end{equation}
where the exponential map of the exact Poincar\'e model is implemented via
\texttt{geoopt}'s \texttt{expmap0} operation. A projection step enforces numerical
stability near the manifold boundary. Decoder layers remain Euclidean, acting on the
embedded hyperbolic samples $\bm{z}$. Although the KL divergence is evaluated using the
Euclidean parameters $(\bm{\mu},\bm{\log\sigma^2})$, the latent codes themselves inhabit
hyperbolic space, facilitating representations aligned with tree-like taxonomic
structures.
\section{Training Loop and Optimization}
Across all model types, the training pipeline \texttt{train\_once}\footnote{\texttt{src/biomevae/trainers/train\_loop.py}}
performs the following steps:
\begin{enumerate}[label=\arabic*., leftmargin=*]
    \item Split the dataset into train/validation partitions using a reproducible random
    shuffle (\texttt{train\_val\_split}).
    \item Optionally standardize features on the training subset and apply the same
    transform to validation and full data.
    \item Instantiate the VAE variant specified by \texttt{model\_type}, configure the
    optimizer (Adam or AdamW), and create data loaders.
    \item Iterate for the requested number of epochs, updating the KL weighting schedule
    and computing losses as described above. Gradient clipping is supported to stabilize
    optimization.
    \item Track training and validation metrics (loss components, latent statistics,
    $R^2$) and trigger early stopping after a patience threshold without validation
    improvement.
    \item Persist the best-performing model checkpoint alongside latent embeddings and
    reconstructions exported as TSV tables.
\end{enumerate}
The same routine underpins Optuna studies: each trial draws hyperparameters, attaches
taxonomy tensors, and invokes \texttt{train\_once}. The best trial is retrained using
its associated seed and serialized for deployment.
\section{Evaluation Utilities}
Beyond training, the repository provides CLI tools for extracting embeddings,
evaluating reconstructions on held-out data, and benchmarking against a
non-negative matrix factorization baseline. All neural models reuse the same
preprocessing stack, ensuring that downstream comparisons remain faithful to the
training configuration. Command-line entry points live under
\texttt{src/biomevae/cli/} and expose flags for selecting models, loss weights, and
evaluation splits. The exported reports include the experiment seed, the
preprocessed data directory, and the path to the checkpoint used for inference so
that results can be reproduced exactly.

\subsection{Latent Space Diagnostics}
The script \texttt{biomevae-embed} materializes low-dimensional embeddings for all
samples in a dataset by loading the serialized configuration and model weights.
Users may request either the encoder means (\texttt{--emb-space mu}) or the
Poincar\'e-ball coordinates for hyperbolic models (\texttt{--emb-space ball}).
The command writes \texttt{embeddings.tsv} to the chosen output directory and,
when \texttt{--export-recon} is supplied, also emits \texttt{recon.tsv} containing
the decoder reconstructions. These tables retain the sample identifiers and are
compatible with downstream ordination or visualization tools.

\subsection{Reconstruction Benchmarks}
Held-out reconstructions are evaluated by \texttt{biomevae-test}, which reuses the
stored preprocessing statistics before scoring the model with the same loss
functions configured during training. The resulting metrics---mean absolute
error, mean squared error, KL divergence statistics, and objective-specific
summaries---are collected in \texttt{test\_report.json}. When taxonomy-aware losses
are active, the routine also aggregates the hierarchical penalties per level to
facilitate comparisons of coarse-grained fidelity across models.

\subsection{Baselines and Comparative Analysis}
The \texttt{biomevae-nmf} entry point computes non-negative matrix factorization
baselines using the identical preprocessing pipeline as the neural models. The
helpers \texttt{biomevae-comparetonmf} and \texttt{biomevae-allcomp} orchestrate
side-by-side evaluations by consuming configuration files produced during
training, yielding JSON summaries that can be appended directly to the LaTeX
tables documenting experimental results.

\section{Experimental Protocol}
The repository includes configuration files under \texttt{configs/} that capture
default hyperparameters for Euclidean, hyperbolic, and taxonomy-regularized VAEs.
Experiments can be launched via the CLI by pointing to one of these YAML
configurations, or by composing equivalent command-line arguments. Each
configuration specifies optimizer settings, learning-rate schedules, KL annealing
strategies, taxonomy weights, and output directories so that individual runs can
be reproduced without modifying the Python source files.

\subsection{Datasets}
We provide example configurations for both synthetic benchmarks and real-world
microbiome cohorts. The synthetic datasets permit controlled studies where the
latent structure is known, while the real-world datasets illustrate behavior on
noisy, sparse counts. Data ingestion scripts normalize feature names, harmonize
taxonomy identifiers, and filter out low-prevalence species to reduce
computational load. Configuration files document the provenance of each table and
its associated taxonomy resource, ensuring that collaborators can reconstruct the
exact training inputs.

\subsection{Reproducibility}
Each training run records a unique identifier combining the dataset name, model
type, and timestamp. Random seeds are set at the level of Python, NumPy, and
PyTorch to minimize nondeterminism. The resulting artifacts---including learned
parameters, latent embeddings, evaluation reports, and reconstructed tables---are
stored in structured directories that can be reloaded by the evaluation
utilities. Versioning these artifacts alongside the LaTeX report preserves a
clear provenance trail and allows automated regeneration of tables and figures.

\section{Conclusion}
The \texttt{biomevae} codebase implements a family of VAEs tailored to microbiome count
matrices, combining hierarchical supervision and non-Euclidean latent geometry. The
framework modularizes data handling, model definition, and regularization, enabling
easy experimentation with taxonomy-aware objectives and hyperbolic representations.
\begin{thebibliography}{9}
\bibitem{kingma2014vae}
D.~P. Kingma and M.~Welling.
\newblock Auto-encoding variational Bayes.
\newblock \emph{International Conference on Learning Representations}, 2014.
\bibitem{higgins2017beta}
I.~Higgins et~al.
\newblock beta-{VAE}: Learning basic visual concepts with a constrained variational framework.
\newblock \emph{International Conference on Learning Representations}, 2017.
\bibitem{nickel2017poincare}
M.~Nickel and D.~Kiela.
\newblock Poincar\'e embeddings for learning hierarchical representations.
\newblock \emph{Advances in Neural Information Processing Systems}, 2017.
\end{thebibliography}
\end{document}
